% Complete Borel--Cantelli, tail and finite-permutation zero-one proofs.
\section{Infinitely often events and zero-one laws: proofs}
\label{sec:annex-p15-limsup-events-and-zero-one-laws}

Use Chapter~\ref{sec:p15-limsup-events-and-zero-one-laws} for the event,
information and permutation definitions. The source proofs are Durrett's
fixed author version \cite[pp.~67--72, 81--83]{DurrettPTEAuthor2019}.

\begin{annexitem}{annex:p15.event-limit-calculus}{thm:p15-event-limit-calculus}{Pointwise event limits and countable convergence thresholds}
\begin{proof}
Both displayed event limits use countably many unions and intersections
of measurable events, so they are measurable. At an outcome \(\omega\),
membership in every \(\bigcup_{n\geq m}A_n\) means that after
every index another occurrence remains, exactly infinitely many
occurrences. Membership in some \(\bigcap_{n\geq m}A_n\) means
membership at every index from that point onward. Negating these
quantifier conditions, or applying De Morgan's laws, gives the two
complement identities.

An indicator sequence has limit superior one exactly when it contains
infinitely many ones; otherwise it is eventually zero. Its limit
inferior is one exactly when it is eventually one. These observations
give the indicator identities and also show that the total count is
infinite exactly on the limit-superior event.

Put \(U_m=\bigcup_{n\geq m}A_n\). Then \(U_m\) decreases
to \(\limsup A_n\), and
\(\mathbb P(U_m)\geq\sup_{n\geq m}\mathbb P(A_n)\).
Continuity from above of the finite probability measure proves
the limit-superior probability inequality. Likewise
\(V_m=\bigcap_{n\geq m}A_n\) increases to \(\liminf A_n\),
and \(\mathbb P(V_m)\leq\inf_{n\geq m}\mathbb P(A_n)\).
Continuity from below proves the other inequality.

If \(X_n\to X\) at an outcome, the deviations exceed any fixed
positive threshold only finitely often. Thus a.s. convergence implies
every stated zero-probability limit-superior condition. Conversely,
remove the union over positive integers \(r\) of the null events
\(\{|X_n-X|>1/r\text{ i.o.}\}\). This is a measurable null
union. At any remaining outcome, for each \(r\) the deviations
are eventually at most \(1/r\). Given \(\varepsilon>0\),
choose \(r\) with \(1/r<\varepsilon\); this proves convergence
at that outcome. The same comparison shows that the integer-threshold
conditions imply the version with every positive real threshold.
No uncountable union of null exceptional sets is taken.
\end{proof}
\end{annexitem}

\begin{annexitem}{annex:p15.borel-cantelli-first}{thm:p15-borel-cantelli-first}{The source counting proof and the tail-union bound}
\begin{proof}
The total count is the increasing limit of its nonnegative measurable
partial sums. The nonnegative series integration law gives
\[
 \mathbb E N_\infty
   =\sum_{n\geq1}\mathbb E\mathbf1_{A_n}
   =\sum_{n\geq1}\mathbb P(A_n).
\]
If this value is finite, the finite-integral criterion makes
\(N_\infty<\infty\) a.s. Since infinitely many occurrences
are exactly \(N_\infty=\infty\), their probability is zero.
This is the source's counting proof, using the previously proved
nonnegative integral theorem.

Countable subadditivity also gives
\(\mathbb P(U_m)\leq\sum_{n\geq m}\mathbb P(A_n)\).
For a convergent nonnegative series, these tails tend to zero,
so continuity from above gives the same conclusion for
\(\limsup A_n=\bigcap_m U_m\). This proves the quantitative
bound without independence. Finally apply the result to the
deviation events separately for each integer \(r\), then use
the countable-threshold convergence criterion just proved.
\end{proof}
\end{annexitem}

\begin{annexitem}{annex:p15.borel-cantelli-second}{thm:p15-borel-cantelli-second}{The mutually independent complement-product proof}
\begin{proof}
Complementing mutually independent events preserves their independence
by Theorem~\ref{thm:p15-independence-basic}. Thus for \(m\leq N\),
\[
 \mathbb P\left(\bigcap_{n=m}^N A_n^c\right)
                      =\prod_{n=m}^N(1-\mathbb P(A_n)).
\]
For \(t\geq0\), the function \(e^{-t}-1+t\) starts at
zero and has derivative \(1-e^{-t}\geq0\), so
\(1-t\leq e^{-t}\). Applying this to each probability in
\([0,1]\) gives the stated exponential upper bound, including
a zero product when one of the probabilities is one.

For fixed \(m\), divergence of the full probability series
implies divergence of \(\sum_{n=m}^N\mathbb P(A_n)\) as
\(N\to\infty\). The upper bound tends to zero. The finite
intersections of complements decrease, so continuity from above
gives \(\mathbb P(\bigcap_{n\geq m}A_n^c)=0\).
Consequently \(\mathbb P(\bigcup_{n\geq m}A_n)=1\) for
every \(m\). Their decreasing intersection is the infinitely-often
event and has probability one. The first lemma handles the
finite-sum alternative; this proof supplies the infinite-sum
alternative under the source's mutual-independence hypothesis.
\end{proof}
\end{annexitem}

\begin{annexitem}{annex:p15.pairwise-event-count-law}{thm:p15-pairwise-event-count-law}{The full pairwise counting refinement and interpolation}
\begin{proof}
The indicators of pairwise independent events are pairwise independent
and hence uncorrelated, with means \(p_j=\mathbb P(A_j)\)
and variances \(p_j(1-p_j)\leq p_j\). Set \(t_0=0\).
The variance-sum theorem
therefore gives
\[
 \mathbb E N_n=t_n,\qquad
 \operatorname{Var}(N_n)=\sum_{j=1}^n p_j(1-p_j)\leq t_n.
\]
For \(t_n>0\), Chebyshev's inequality yields
\[
 \mathbb P(|N_n-t_n|>\delta t_n)
                  \leq\frac1{\delta^2t_n}\qquad(\delta>0).
\]
This already proves convergence of the ratio in probability when
\(t_n\to\infty\), but the source obtains a.s. convergence
by the following additional argument.

Let \(n_k=\min\{n:t_n\geq k^2\}\), which exists for
every \(k\geq1\). Since each increment of \(t_n\) is at
most one, minimality gives
\(k^2\leq t_{n_k}\leq k^2+1\).
The indices strictly increase, because \(k^2+1<(k+1)^2\),
and tend to infinity since \(t_n\leq n\) implies \(n_k\geq k^2\).
Put \(T_k=N_{n_k}\). The displayed deviation bound gives
\[
 \mathbb P(|T_k-t_{n_k}|>\delta t_{n_k})
                         \leq\frac1{\delta^2k^2}.
\]
The series of these bounds converges: for \(k\geq2\),
\(1/k^2\leq1/(k-1)-1/k\), so the tail is bounded by a
telescoping series. The first Borel--Cantelli lemma, applied
for each \(\delta=1/r\), \(r\geq1\), shows on one
probability-one event that \(T_k/t_{n_k}\to1\).

At an outcome in that event and for \(n_k\leq n<n_{k+1}\),
both the occurrence counts and their means are nondecreasing, so
\[
 \frac{T_k}{t_{n_{k+1}}}
             \leq\frac{N_n}{t_n}
             \leq\frac{T_{k+1}}{t_{n_k}}.
\]
All denominators are positive here. The lower bound is
\((T_k/t_{n_k})(t_{n_k}/t_{n_{k+1}})\), and the upper
bound is \((T_{k+1}/t_{n_{k+1}})(t_{n_{k+1}}/t_{n_k})\).
The remaining deterministic ratio satisfies
\[
 1\leq\frac{t_{n_{k+1}}}{t_{n_k}}
              \leq\frac{(k+1)^2+1}{k^2}\longrightarrow1.
\]
Thus both bounds tend to one. Every sufficiently large \(n\)
lies between successive \(n_k\), with \(k\to\infty\),
so the full ratio tends to one a.s. Since \(t_n\to\infty\),
the counts \(N_n\) are unbounded at those outcomes. This is
exactly infinitely many occurrences. No mutual-independence
step was used in this proof.
\end{proof}
\end{annexitem}

\begin{annexitem}{annex:p15.finite-coordinate-approximation}{thm:p15-finite-coordinate-approximation}{Approximation by the increasing finite-information fields}
\begin{proof}
Let \(\mathcal C\) contain those events \(A\in\mathcal H_\infty\)
that admit the asserted approximation for every positive error.
Every event in \(\bigcup_n\mathcal H_n\) belongs to
\(\mathcal C\), using itself as its approximation.
Complementation preserves errors because
\(A^c\triangle B^c=A\triangle B\), so \(\mathcal C\)
is closed under complements.

For \(A_j\in\mathcal C\), set \(A=\bigcup_{j\geq1}A_j\).
Given \(\varepsilon>0\), continuity from below permits a
positive integer \(M\) with
\(\mathbb P(A\setminus\bigcup_{j=1}^M A_j)<\varepsilon/2\).
For each \(j\leq M\), choose \(B_j\in\mathcal H_{m_j}\)
with \(\mathbb P(A_j\triangle B_j)<\varepsilon/(2M)\).
The finite union \(B=\bigcup_{j=1}^M B_j\) belongs to
\(\mathcal H_{\max_jm_j}\), and
\[
 A\triangle B\subseteq
    \left(A\setminus\bigcup_{j=1}^M A_j\right)
       \cup\bigcup_{j=1}^M(A_j\triangle B_j).
\]
Its probability is less than \(\varepsilon\). Hence
\(\mathcal C\) is a sigma-algebra containing the generators
of \(\mathcal H_\infty\), proving the approximation claim.

For a full sequence with the stated coordinate index, let
\(d_n=\inf_{B\in\mathcal H_n}\mathbb P(A\triangle B)\).
The infimum is over a nonempty set and lies in \([0,1]\).
It decreases to zero by the just-proved approximation. Choose
\(A_n\in\mathcal H_n\) with
\(\mathbb P(A_n\triangle A)<d_n+1/n\). Such a choice exists
by the definition of the infimum, and its error tends to zero.

The probability difference identity
\(\mathbb P(A)-\mathbb P(B)=\mathbb P(A\setminus B)
-\mathbb P(B\setminus A)\) gives the absolute bound in
the statement. If membership in \(A\) and \(C\) differs,
it must differ between \(A,B\) or between \(B,C\), proving
the symmetric-difference inclusion.

On the canonical product, let \(\pi_n\) be the first-\(n\)
coordinate map. The collection
\(\{\pi_n^{-1}(B):B\in\mathcal S^{\otimes n}\}\) is a
sigma-algebra containing each of the first \(n\) coordinate
sigma-algebras. Conversely \(\pi_n\) is measurable for
\(\mathcal H_n\), since each rectangle inverse image is
a finite intersection of its coordinate events. The collection
therefore equals \(\mathcal H_n\). This proves the exact
finite-coordinate representation used in the source's block-swap
argument, with no appeal to a completion or conditional expectation.
\end{proof}
\end{annexitem}

\begin{annexitem}{annex:p15.tail-event-structure}{thm:p15-tail-event-structure}{Remote-future measurability and canonical finite changes}
\begin{proof}
An intersection of sigma-algebras on the same space is a sigma-algebra,
so \(\mathcal T\) is one. For any fixed \(m\), deleting
finitely many initial events leaves their limit superior unchanged.
All events \(\{X_n\in B_n\}\) with \(n\geq m\) belong
to \(\mathcal F'_m\). Their countable union/intersection
description therefore puts the infinitely-often event in that
field for every \(m\), hence in \(\mathcal T\).

For real coordinates, put \(S_0=0\). At a fixed cutoff \(m\),
the variables \(R_n^{(m)}=\sum_{i=m}^nX_i\), \(n\geq m\),
are \(\mathcal F'_m\)-measurable, and
\(S_n=R_n^{(m)}+S_{m-1}\). Adding that finite real constant
at each outcome preserves convergence to a finite real limit,
and preserves convergence to either infinite endpoint.
The event of finite convergence of \(R_n^{(m)}\) is measurable
by the Cauchy criterion
\[
 \bigcap_{r\geq1}\bigcup_{N\geq m}
       \bigcap_{j,k\geq N}\{|R_j^{(m)}-R_k^{(m)}|\leq1/r\}.
\]
The real numbers are complete, so this criterion is equivalent
to finite convergence. Convergence to \(+\infty\) or
\(-\infty\) likewise has a countable eventual-threshold
description, using integer thresholds. These events belong to
every \(\mathcal F'_m\), as asserted.

If \(c_n\to+\infty\), then for all sufficiently large \(n\)
the ratios are defined and
\(S_n/c_n-R_n^{(m)}/c_n=S_{m-1}/c_n\to0\) pointwise.
Two real sequences with difference tending to zero have the same
extended limit superior: eventually each is at most the other
plus any given \(\varepsilon>0\), giving both limsup
inequalities, and then let \(\varepsilon\downarrow0\).
The extended limit superior of the \(\mathcal F'_m\)-measurable
ratios is measurable by countable suprema and infima. Thus its
strict and weak threshold events, identical to those for the
original ratios, belong to every future field.

For the canonical criterion, fix \(m\). Events whose membership
depends only on coordinates \(i\geq m\) form a sigma-algebra
containing the future-coordinate generators. Every event of
\(\mathcal F'_m\) therefore has that property. A tail event
is consequently unchanged by every modification of finitely many
coordinates, by choosing a cutoff past all the modified indices.

Conversely, take a product-measurable event \(A\) with that
invariance, and choose one \(s_0\in S\). Define \(r_m\)
to replace the first \(m-1\) coordinates by \(s_0\) and
leave all later coordinates unchanged. As a map into the product
space, \(r_m\) is measurable with respect to \(\mathcal F'_m\):
its coordinate maps are constants or future coordinates, and the
target sigma-algebra is generated by those coordinates. Invariance
gives \(A=r_m^{-1}(A)\), so \(A\in\mathcal F'_m\) for
every \(m\). This proves the converse exactly, not merely a.s.

On \(\mathbb R^{\mathbb N_+}\), the sequences
\((1,0,0,\ldots)\) and \((-1,0,0,\ldots)\) differ only
at their first coordinate. Their partial sums have limit superiors
one and minus one. The event \(\{\limsup S_n>0\}\) distinguishes
them, so it is not tail by the necessary criterion. This nonmembership
example is a statement on that canonical space; it is not a universal
nonmembership assertion about every degenerate or dependent process.
\end{proof}
\end{annexitem}

\begin{annexitem}{annex:p15.kolmogorov-tail-zero-one}{thm:p15-kolmogorov-tail-zero-one}{The source's two independence extensions and the self-independence step}
\begin{proof}
Fix \(k\geq1\). For every \(j\geq1\), mutual independence
and the finite disjoint-group theorem make \(\mathcal H_k\)
independent of \(\sigma(X_{k+1},\ldots,X_{k+j})\).
The union over \(j\) of these latter, increasing sigma-algebras
is a pi-system containing \(\Omega\). It is independent of
\(\mathcal H_k\), since any selected event lies in one of
the finite future blocks. The independent-generator theorem
extends this to independence of \(\mathcal H_k\) and
\(\mathcal F'_{k+1}\). This is the source's first step.

The union \(\mathcal C=\bigcup_k\mathcal H_k\) is also a
pi-system containing \(\Omega\): any finite intersection
lies in the largest of its finitely many index fields.
For \(B\in\mathcal C\), choose \(k\) with \(B\in\mathcal H_k\).
For \(A\in\mathcal T\), the inclusion
\(\mathcal T\subseteq\mathcal F'_{k+1}\) and the first step
give \(\mathbb P(A\cap B)=\mathbb P(A)\mathbb P(B)\).
Thus the pi-systems \(\mathcal C\) and \(\mathcal T\)
are independent. Applying the generator theorem again gives
independence of \(\mathcal H_\infty=\sigma(\mathcal C)\)
and \(\mathcal T\), the source's second step.

Every \(A\in\mathcal T\) is also in \(\mathcal H_\infty\).
It is therefore independent of itself, and
\(\mathbb P(A)=\mathbb P(A\cap A)=\mathbb P(A)^2\).
The only values in \([0,1]\) satisfying this equality are
zero and one. These arguments use mutual independence but
neither identical laws nor integrability assumptions.
\end{proof}
\end{annexitem}

\begin{annexitem}{annex:p15.permutation-structure}{thm:p15-permutation-structure}{Exact event invariance and invariance of an i.i.d. law}
\begin{proof}
Each coordinate of \(T_\pi\) is a coordinate projection, so
the map is measurable for the product sigma-algebra. Its inverse
is \(T_{\pi^{-1}}\), which is measurable by the same argument.
The whole space is invariant. Inverse images commute with
complements and countable unions, so these operations preserve
invariance under every finite permutation. Thus \(\mathcal E\)
is a sigma-algebra. A finite permutation changes only finitely
many coordinates of any one sequence. The canonical finite-change
criterion consequently gives \(\mathcal T\subseteq\mathcal E\).

Suppose the coordinate variables are i.i.d. with common law
\(\eta\). For any finite cylinder specified by measurable
value sets at distinct indices, its inverse image under \(T_\pi\)
uses the same value sets at the permuted distinct indices.
Mutual independence and identical laws give the same probability,
the product of the corresponding \(\eta\)-values. Thus
\(\mathbb P\circ T_\pi^{-1}\) and \(\mathbb P\) are
probabilities agreeing on the cylinder pi-system that generates
the product sigma-algebra. The uniqueness theorem makes them
equal on that entire sigma-algebra. This proves law invariance.

For real coordinates, choose \(m\) past the finite support
of \(\pi\). The permutation maps \(\{1,\ldots,m\}\)
onto itself and fixes all later indices. For \(n\geq m\),
the sum \(S_n(T_\pi\omega)\) therefore equals \(S_n(\omega)\).
Discarding finitely many terms changes neither infinitely-often
membership nor a limit superior. This proves both source families
of permutable events, with ratios interpreted only where their
deterministic denominators are eventually nonzero.

For strict inclusion, on the canonical binary product let
\(A=\bigcup_{n\geq1}\{X_n=1\}\). Permuting coordinates
leaves this event unchanged. But it distinguishes the paths
\((1,0,0,\ldots)\) and \((0,0,0,\ldots)\), so it is not
tail. This is a property of the measurable sets themselves and
persists under any probability law on that product, including
the already constructed i.i.d. fair-bit law. Hence probability
triviality must not be substituted for literal equality of
\(\mathcal E\) and \(\mathcal T\).
\end{proof}
\end{annexitem}

\begin{annexitem}{annex:p15.hewitt-savage-zero-one}{thm:p15-hewitt-savage-zero-one}{Finite approximation and the independent swapped block}
\begin{proof}
Let \(A\in\mathcal E\). On this canonical product, the
sigma-algebra generated by all coordinates is the entire stated
product sigma-algebra. The finite-information approximation theorem
gives \(A_n\in\mathcal H_n\) with
\(\varepsilon_n=\mathbb P(A_n\triangle A)\to0\), and
\(A_n=\{(X_1,\ldots,X_n)\in B_n\}\) for some
\(B_n\in\mathcal S^{\otimes n}\).

Define the finite block-swap permutation by
\[
 \pi^{(n)}(j)=\begin{cases}
 j+n,&1\leq j\leq n,\\
 j-n,&n<j\leq2n,\\
 j,&j>2n.
 \end{cases}
\]
It is its own inverse. Set
\[
 A_n'=T_{\pi^{(n)}}^{-1}(A_n)
        =\{(X_{n+1},\ldots,X_{2n})\in B_n\}.
\]
The events \(A_n,A_n'\) belong to disjoint finite coordinate
blocks and are independent. The law is invariant under this
permutation by the preceding theorem, and \(A\) is exactly
invariant by assumption. Inverse images preserve symmetric
difference, so
\[
 \mathbb P(A_n'\triangle A)
   =\mathbb P(T_{\pi^{(n)}}^{-1}(A_n\triangle A))
   =\mathbb P(A_n\triangle A)=\varepsilon_n.
\]
The probability-difference bound therefore gives
\(\mathbb P(A_n)\to\mathbb P(A)\) and
\(\mathbb P(A_n')\to\mathbb P(A)\).
The symmetric-difference triangle inclusion also gives
\[
 \mathbb P(A_n\triangle A_n')\leq2\varepsilon_n\to0.
\]
In particular,
\[
 0\leq\mathbb P(A_n)-\mathbb P(A_n\cap A_n')
     =\mathbb P(A_n\setminus A_n')
     \leq\mathbb P(A_n\triangle A_n')\to0.
\]
It follows that \(\mathbb P(A_n\cap A_n')\to\mathbb P(A)\).
But independence of the two blocks gives
\(\mathbb P(A_n\cap A_n')=\mathbb P(A_n)\mathbb P(A_n')
\to\mathbb P(A)^2\). Thus \(\mathbb P(A)=\mathbb P(A)^2\),
and its value is zero or one. This supplies all approximation,
permutation and limiting steps of the source's proof.

For an i.i.d. sequence on another space, its map into the
canonical coordinate product is measurable, since cylinder
inverse images are finite intersections of measurable value
events. Its pushforward law makes the canonical coordinates
i.i.d. with the same finite-dimensional laws. Apply the result
just proved there and use the definition of pushforward to
obtain the pullback conclusion.
\end{proof}
\end{annexitem}

\begin{annexitem}{annex:p15.zero-one-measurable-values}{thm:p15-zero-one-measurable-values}{The rational-cut argument for an almost-sure constant}
\begin{proof}
For every rational \(q\), measurability gives
\(p(q)=\mathbb P(Z\leq q)\in\{0,1\}\).
The function \(p\) is nondecreasing. Let
\(U=\{q\in\mathbb Q:p(q)=1\}\).
If \(U\) is empty, the countable union of events
\(\{Z\leq n\}\), \(n\geq1\), has probability zero.
That union is \(\{Z<\infty\}\), so \(Z=\infty\) a.s.
If \(U=\mathbb Q\), every event \(\{Z>-n\}\) is null;
their union is \(\{Z>-\infty\}\), so \(Z=-\infty\) a.s.

Otherwise \(U\) is nonempty and has a finite lower bound:
any rational \(q_0\notin U\) has \(p(q_0)=0\), and by
monotonicity no smaller rational belongs to \(U\).
Its infimum \(c\) is a finite real number. For rational
\(q<c\), \(p(q)=0\). For rational \(q>c\), some member
of \(U\) lies below \(q\), so \(p(q)=1\).
Choose rational lower and upper bounds approaching \(c\)
from the two sides. Each event that \(Z\) is above the
lower bound and at most the upper bound has probability one.
Their countable intersection forces \(Z=c\), proving the
remaining case. If \(Z\) is finite a.s., the two infinite
constant cases are excluded. Applying this general result to
the sigma-algebras proved probability-trivial above gives the
stated applications, without claiming a value for the constant
or existence of an unrelated sequence limit.
\end{proof}
\end{annexitem}

\begin{annexitem}{annex:p15.zero-one-countermodels}{thm:p15-zero-one-countermodels}{The explicit probability models and all failed hypotheses}
\begin{proof}
For any sequence \(p_n\in[0,1]\), the countable product
construction applied to the two-point laws with masses
\(p_n\) at one and \(1-p_n\) at zero supplies mutually
independent binary coordinates. For \(p_n=1/n\), the
harmonic series diverges: the block of indices from
\(2^k+1\) through \(2^{k+1}\) contributes at least one half.
The second Borel--Cantelli lemma gives infinitely many ones a.s.,
and the pairwise counting theorem gives \(N_n/H_n\to1\)
a.s. For \(p_n=1/n^2\), the telescoping comparison used
above proves summability, and the first lemma gives only
finitely many ones a.s.

On the uniform space, \(\mathbb P((0,1/n))=1/n\).
For each outcome \(u>0\), every sufficiently large \(n\)
satisfies \(1/n<u\), so that outcome belongs to none of
the later intervals. Their limit superior is therefore empty,
although the probability sum diverges. More generally the
same pointwise argument holds for \(A_n=(0,a_n)\),
\(0\leq a_n\leq1\), whenever \(a_n\to0\); a lower
bound \(a_n\geq1/n\) makes the probability series divergent.
In the particular case \(a_n=1/n\), the events are nested and are not independent:
\(\mathbb P(A_2\cap A_3)=1/3\ne(1/2)(1/3)\).

For the pairwise-independent countermodel, the required independent
fair bits \(Y,B_1,B_2,\ldots\) exist by the earlier bit
construction and reindexing. For outputs in the same pair,
\((B_j,Y\oplus B_j)=(u,v)\) fixes exactly
\(B_j=u\), \(Y=u\oplus v\), a joint outcome of probability
one quarter. Thus those two coordinates are independent fair bits.
For coordinates in different pairs, write them as
\(B_j\oplus aY\) and \(B_k\oplus bY\), where
\(j\ne k\) and \(a,b\in\{0,1\}\). For each prescribed
output pair \((u,v)\), either value of \(Y\) determines
exactly one pair of values of \(B_j,B_k\). There are two
of the eight equally likely outcomes of the independent triple
\((Y,B_j,B_k)\), giving probability one quarter again.
This verifies pairwise independence and identical fair marginals
for every pair of output coordinates.

At every outcome, \(X_{2j-1}\oplus X_{2j}=Y\). Given any
tail cutoff \(m\), choose \(j\) with \(2j-1\geq m\).
Then \(Y\) is a measurable function of two coordinates
in that future field, so \(\{Y=1\}\in\mathcal F'_m\).
It belongs to every such field and has probability one half.
The family is not mutually independent: the finite output
pattern \((X_1,X_2,X_3,X_4)=(0,0,0,1)\) is impossible,
since its two pair parities differ, whereas four independent
fair bits would assign it probability \(1/16\).

For the canonical event \(A=\{\text{some coordinate is one}\}\),
permutability and failure of tail membership were proved in
the permutation-structure argument using the one-hot and zero
paths. Under the i.i.d. fair-bit law,
\[
 \mathbb P(A^c)=\lim_{n\to\infty}
      \mathbb P(X_1=\cdots=X_n=0)=\lim_{n\to\infty}2^{-n}=0.
\]
Under the independent law with only the first coordinate random,
each later event \(\{X_n=1\}\) is null. Their countable
union is null, so \(A\) agrees a.s. with \(\{X_1=1\}\)
and has probability one half. The coordinates are independent
by their product construction, but their marginal laws differ.
This refutes removal of identical distributions from Hewitt--Savage.

Finally construct a fair bit \(Y\) on a two-point space and
map it to the constant coordinate sequence \((Y,Y,\ldots)\).
This is a measurable map into the canonical binary product,
so its pushforward is a probability law. It gives one half
of the mass to the all-zero path and one half to the all-one
path. Both paths are fixed by every coordinate permutation,
so the law is invariant under every finite permutation. All
marginal laws are fair binary, but two-coordinate independence
already fails: \(\mathbb P(X_1=1,X_2=1)=1/2\ne1/4\).
The event of infinitely many ones is tail and permutable,
contains exactly the all-one path among these two support
points, and therefore has probability one half. This proves
the stated failure despite exchangeability of the law.
\end{proof}
\end{annexitem}
